\documentclass{beamer}
\usepackage{graphicx}
\usepackage{xcolor}
\usepackage{tikz}
\usepackage{amsmath}

\definecolor{purpleblue}{RGB}{80, 80, 180}
\definecolor{novelPurple}{RGB}{98, 54, 255}
\definecolor{softPurple}{RGB}{180, 170, 255}
\definecolor{softGray}{RGB}{235, 235, 245}
\definecolor{lightText}{RGB}{90, 90, 110}

\usetheme{Madrid}
\usecolortheme{dolphin}
\graphicspath{{static/course_materials/}{static/hard/}{./}}

\setbeamerfont{title}{size=\Large,series=\bfseries}
\setbeamerfont{frametitle}{size=\LARGE,series=\bfseries}
\setbeamercolor{frametitle}{fg=white, bg=novelPurple}
\setbeamercolor{title}{fg=white, bg=purpleblue}

\title[SAT Hard Question 1]{\textbf{SAT Hard Question\\ Set 1}}
\author{\textbf{Jack Zeng}}
\institute{\textcolor{white}{\textbf{Novel Prep}}}
\date{\today}

\begin{document}

\begin{frame}
    \titlepage
    \vfill
    \begin{flushright}
        \textcolor{purpleblue}{\Huge \textbf{Novel Prep}}
    \end{flushright}
\end{frame}

\begin{frame}
    \begin{tikzpicture}[remember picture, overlay]
        \shade[top color=softPurple, bottom color=softGray]
            (current page.north west) rectangle (current page.south east);
        \node[align=center] at (current page.center) {
            {\Huge\bfseries\textcolor{purpleblue}{Novel Prep}} \\[0.5cm]
            {\Large\itshape\textcolor{lightText}{Phase 2 · Hard Question Practice}}
        };
    \end{tikzpicture}
\end{frame}

\begin{frame}{Overview}
    \tableofcontents
\end{frame}

\begin{frame}
\section{Hard Question Set 1}
    \begin{tikzpicture}[remember picture, overlay]
        \shade[top color=softPurple, bottom color=softGray]
            (current page.north west) rectangle (current page.south east);
        \node[align=center] at (current page.center) {
            {\LARGE\bfseries\textcolor{purpleblue}{Hard Question Set 1}} \\[0.5cm]
            {\Large\itshape\textcolor{lightText}{11 SAT-style challenge problems}}
        };
    \end{tikzpicture}
\end{frame}

% --- Question 1 ---
\begin{frame}{Question 1}
\small
In the given expression, \(a\) and \(c\) are positive constants. If \(px + q\) is a factor of the expression,
\[
ax^2 + 176x + c = (px + q)(rx + s),
\]
where \(p\) and \(q\) are positive constants, what is the greatest possible value of \(ac\)?
\end{frame}
\begin{frame}{Answer 1}
\small
\textbf{Final Answer:} $\boxed{7744}$
\end{frame}

% --- Question 2 ---
\begin{frame}{Question 2}
\small
In the given expression,
\[
34z^{18} + b z^9 + 30,
\]
\(b\) is a positive integer. If \(qz^9 + r\) is a factor of the expression, where \(q\) and \(r\) are positive integers, what is the greatest possible value of \(b\)?
\end{frame}
\begin{frame}{Answer 2}
\small
\textbf{Final Answer:} $\boxed{1021}$
\end{frame}

% --- Question 3 ---
\begin{frame}{Question 3}
\small
The expression
\[
6x^4 + 17x^2 + 7
\]
can be rewritten as \((3x^2 + a)(2x^2 + b)\), where \(a\) and \(b\) are positive integers, or as \((3x^2 + c)(2x^2 + d)\), where \(c\) and \(d\) are positive non-integers.

What is the value of \(a + c\)?
\end{frame}
\begin{frame}{Answer 3}
\small
\textbf{Final Answer:} $\boxed{\dfrac{17}{2}}$
\end{frame}

% --- Question 4 ---
\begin{frame}{Question 4}
\small
The given quadratic function
\[
54x^4 + 219x^2 + 105
\]
has factors in the form \((k)(ax^2 + b)(cx^2 + d)\). If \(a, b, c, d, k\) are integers, what is the smallest possible value of \(ab\)?
\end{frame}
\begin{frame}{Answer 4}
\small
\textbf{Final Answer:} $\boxed{14}$
\end{frame}

% --- Question 5 ---
\begin{frame}{Question 5}
\small
The equation of a parabola is written in the form
\[
y = ax^2 + bx + c,
\]
where \(a, b, c\) are constants. When graphed in the \(xy\)-plane, the parabola has vertex \((-1, 4)\) and does not intersect the \(x\)-axis. Which of the following could be the value of \(a - b - c\)?
\vspace{0.35cm}
\begin{enumerate}
    \item[A.] 2
    \item[B.] 4
    \item[C.] -1
    \item[D.] -6
\end{enumerate}
\end{frame}
\begin{frame}{Answer 5}
\small
\textbf{Correct Answer:} D
\end{frame}

% --- Question 6 ---
\begin{frame}{Question 6}
\small
A quadratic function with vertex \((1, 32)\) can be written as
\[
 -ax^2 + bx + c \quad (a, b, c \text{ are all positive integers}).
\]
What is the maximum value of \(b\)?
\end{frame}
\begin{frame}{Answer 6}
\small
\textbf{Final Answer:} $\boxed{62}$
\end{frame}

% --- Question 7 ---
\begin{frame}{Question 7}
\small
An object's speed is increasing at a rate of \(7.7\ \text{meters per second squared}\). What is this rate in \textbf{miles per minute squared}, rounded to the nearest tenth? (Use \(1\ \text{mile} = 1609\ \text{meters}\).)
\vspace{0.35cm}
\begin{enumerate}
    \item[A.] 0.3
    \item[B.] 17.2
    \item[C.] 206.5
    \item[D.] 209
\end{enumerate}
\end{frame}
\begin{frame}{Answer 7}
\small
\textbf{Correct Answer:} B
\end{frame}

% --- Question 8 ---
\begin{frame}{Question 8}
\small
The function \(f\) is defined by
\[
f(x) = ab^{\frac{x}{n}},
\]
where \(a, b, n\) are constants, and \(b, n\) are integers. If \(f(2) = 6\) and \(f(4) = 150\), what is the value of \(f(5)\)?
\end{frame}
\begin{frame}{Answer 8}
\small
\textbf{Final Answer:} $\boxed{750}$
\end{frame}

% --- Question 9 ---
\begin{frame}{Question 9}
\small
The graph of \(y = f(x)\) in the \(xy\)-plane has a vertex at \((4, -7)\), contains the point \((3, -9)\), and has a \(y\)-intercept at \((0, a)\). The graph of \(y = 6 \cdot f(x)\) has a \(y\)-intercept at \((0, b)\). What is the positive difference between \(a\) and \(b\)?
\end{frame}
\begin{frame}{Answer 9}
\small
\textbf{Final Answer:} $\boxed{195}$
\end{frame}

% --- Question 10 ---
\begin{frame}{Question 10}
\small
The function \(f\) is defined by
\[
f(x) = ax^2 + bx + c,
\]
where \(a, b, c\) are constants. The graph of \(y = f(x)\) passes through the points \((9, 0)\) and \((-2, 0)\). If \(a\) is an integer greater than 1, which of the following could be the value of \(a + b\)?
\vspace{0.35cm}
\begin{enumerate}
    \item[A.] 8
    \item[B.] 7
    \item[C.] -6
    \item[D.] -12
\end{enumerate}
\end{frame}
\begin{frame}{Answer 10}
\small
\textbf{Correct Answer:} D
\end{frame}

% --- Question 11 ---
\begin{frame}{Question 11}
\small
Which of the following expressions has a factor of \(x + 2b\), where \(b\) is a positive integer constant?
\vspace{0.35cm}
\begin{enumerate}
    \item[A.] \(3x^2 + 7x + 14b\)
    \item[B.] \(3x^2 + 22x + 14b\)
    \item[C.] \(3x^2 + 30x + 14b\)
    \item[D.] \(3x^2 + 37x + 14b\)
\end{enumerate}
\end{frame}
\begin{frame}{Answer 11}
\small
\textbf{Correct Answer:} D
\end{frame}

\begin{frame}{}
    \begin{tikzpicture}[remember picture, overlay]
        \shade[top color=novelPurple!40, bottom color=white]
            (current page.north west) rectangle (current page.south east);
        \fill[white, opacity=0.15] (current page.center) circle (5cm);
        \foreach \x in {1,2,3,4,5} {
            \fill[novelPurple, opacity=0.12] (rand*10-5, rand*7-3.5) circle (0.1+\x*0.05);
        }
    \end{tikzpicture}
    \centering
    \vspace{5em}
    {\Huge \textbf{\textcolor{novelPurple}{Thank You!}}} \\[1em]
    {\large \textcolor{gray}{We appreciate your time and focus.}} \\[3em]
    {\LARGE \textbf{\textcolor{novelPurple}{\textsf{Novel Prep}}}}
    \vfill
\end{frame}

\end{document}
